% ===================================================================
%  नकसा नं. १२ — इस्टेसन। --- रेल। --- जहाज।
%
%  Traduction, faite en 2026, en bhojpouri d'aujourd'hui, sur l'ido
%  de texto/io. Regles de la colonne : voir 10-nakasa-01.tex. Une
%  seule numerotation. 28 blocs, 87 renvois, vingt-sept paires a
%  retourner. UNE note d'auteur, t12-noto-1, posee au milieu de
%  l'alinea 14 ; la traduction ecrit UN bloc et pose la note apres.
%  Le « Noto. » d'ouverture n'a pas de cle propre : il fait partie de
%  l'apparat, comme dans l'ido.
%
%  « इस्टेसन » FERME LE CERCLE OUVERT AU TABLEAU 1. L'en-tete du
%  premier tableau posait que le bhojpouri ne supporte pas qu'un mot
%  commence par s + consonne et met une voyelle devant : स्कूल devient
%  इसकूल, et le titre du tableau 1 etait deja par la une preuve de
%  langue. Le voici sur le mot le plus entendu de la region : स्टेशन
%  devient इस्टेसन. LE MEME PROCEDE, LE MEME PRETEUR, ONZE TABLEAUX
%  PLUS TARD — sauf que celui-ci n'est pas un mot d'ecole, c'est le
%  mot d'un lieu que ce pays-la connait mieux qu'aucun autre.
%
%  CAR CE TABLEAU EST, DE TOUT LE LIVRET, CELUI QUI TOMBE LE PLUS
%  JUSTE. Une salle des pas perdus pleine de voyageurs, un homme qui
%  fait peser sa malle, un porteur, un train de nuit : ce n'est pas
%  une scene exotique dans le Bhojpour, c'est sa matiere. C'est de
%  ces gares-la — Chapra, Ballia, Gorakhpur, Buxar — que sont partis
%  les engages du girmit vers Maurice, la Guyane, Trinidad et les
%  Fidji entre 1834 et 1917, et c'est de la meme quai que partent
%  aujourd'hui les migrants de Delhi, du Pendjab et de Bombay. La
%  moitie de la chanson populaire bhojpourie se passe sur un quai. Le
%  tableau 7 avait trouve une ferme qui ne demandait aucun refus ;
%  celui-ci trouve une gare qui n'en demande pas davantage, et pour
%  une raison toute differente — non parce que la chose est ancienne
%  et partagee, mais parce qu'elle est recente et centrale.
%
%  ET « डिब्बा » LE WAGON (37) EST LE SECOND MOT DE FABRICATION
%  LOCALE. Le tableau 10 avait releve पनडुब्बी, le sous-marin, seul
%  mot de mer que la plaine n'ait emprunte a personne, parce qu'elle
%  avait deja le verbe. Ici, sur un vocabulaire ferroviaire
%  entierement anglais — रेल, इस्टेसन, टिकट, प्लेटफार्म, गार्ड,
%  सिगनल, इंजन —, un mot est du fonds : डिब्बा, la boite. Le wagon a
%  ete nomme par ce a quoi il ressemblait, et le compartiment (38)
%  s'appelle खाना, la case. ON N'EMPRUNTE PAS CE QU'ON PEUT DIRE AVEC
%  CE QU'ON A : deuxieme verification de la meme loi, deux tableaux
%  apres la premiere.
%
%  Enfin le porteur (16) est « कुली », et le mot est sorti d'ici. La
%  colonne gujaratie l'avait releve comme l'un des deux seuls mots
%  partis de l'Inde vers l'Europe au lieu d'y entrer — coolie en
%  anglais, coolie en francais, Kuli en allemand. Il est parti par ces
%  gares-la, avec les hommes.
% ===================================================================

%%K t12-apar-1 apar
\VUsaut{4.0mm}
\VUtitre{15.0pt}{60}{नकसा नं. १२}
\VUsaut{2.4mm}
\VUpk{11.4pt}{40}{इस्टेसन। --- रेल। --- जहाज।}
\VUsaut{3.4mm}
\textsc{नोट।} --- \textit{हर देस में रेल के समय-सारनी अलग होले, एहीसे
हमनी के मिसाल खातिर} \VUgras{फ्रेंच सारनी} \textit{के समय लेत बानी जा।}

%%K t12-01-1 p
१. --- हम अभिए जर्मनी के जातरा कइनी। निकले से पहिले हम एगो \VUgras{रेल
के समय-पोथी} \textsuperscript{(15)} खरीदनी ताकि गाड़ी छूटे के समय जान
सकीं। ई देख के कि सबसे सुभीता के गाड़ी पेरिस से राति नौ बजे छूटेले आ
स्ट्रासबुर्ग एक बजके तेरह मिनट पर पहुँचेले, हम आपन जातरा के तइयारी
कइनी, आ अगिले दिने हम आपन के इस्टेसन ले जाए के कहनी।

%%K t12-02-1 p
२. --- जब हम बड़का रवानगी-हाल में घुसनी, एगो बुढ़उ देहाती
\VUgras{पादरी} \textsuperscript{(2)} के पाछे, जे हाथ में आपन \VUgras{जातरा
के झोरा} \textsuperscript{(3)} लटकवले रहलें, बहुते \VUgras{जातरी}
\textsuperscript{(1)} पहिलही से आ चुकल रहलें।

%%K t12-02-2 p
काहेकि हमरा एगो \VUgras{तार} \textsuperscript{(5)} भेजे के रहे, हम एगो
\VUgras{जाँचकर्ता} \textsuperscript{(6)} से \VUgras{तारघर}
\textsuperscript{(4)} पूछ लेनी।

%%K t12-03-1 p
३. --- \VUgras{अगोरे के कमरा} \textsuperscript{(7)} में जाए से पहिले हम
लवटे के \VUgras{टिकट} \textsuperscript{(9)} ले लेनी।

%%K t12-04-1 p
४. --- हमरा \VUgras{खिड़की} \textsuperscript{(8)} पर देर तक ना अगोरे के
पड़ल, काहेकि ओहिजा थोड़े लोग रहे। एगो \VUgras{फौजी} \textsuperscript{(10)},
जे छुट्टी पर जात रहे, भलमनसाहत से आपन नंबर हमरा के दे देलस, एहसे हमरा
लगे एतना समय रह गइल कि हम \VUgras{इस्टेसन मास्टर} \textsuperscript{(13)}
से कुछ बात पूछ सकीं, जे निगरानी के \VUgras{दरोगा} \textsuperscript{(12)} आ
डिउटी पर के \VUgras{पुलिस के सिपाही} \textsuperscript{(11)} से बतियावत
रहलें।

%%K t12-tit-1 sub
\VUpk{10.2pt}{40}{असबाब के दर्ज।}
\VUsaut{3.0mm}

%%K t12-05-1 p
५. --- \VUgras{किताब के गुमटी} \textsuperscript{(14)} से अखबार खरीद के
हम आपन \VUgras{असबाब} \textsuperscript{(17)} तउलवनी, जेकरा के एगो
\VUgras{कुली} \textsuperscript{(16)} अभिए \VUgras{असबाब के ठेला}
\textsuperscript{(19)} पर ले आइल रहे ; \VUgras{बाबू}
\textsuperscript{(22)} जे \VUgras{तराजू} \textsuperscript{(21)} पर लागल
रहलें, हमरा
बतवलें कि हमार संदूक पैंतीस किलो के बा ; एहसे पाँच किलो \VUgras{जादे
वजन} भइल, जेकर खातिर हमरा \VUgras{असबाब के खिड़की}
\textsuperscript{(23)} पर बेसी पइसा देवे के पड़ल।

%%K t12-06-1 p
६. --- काहेकि हम बहुते जल्दी आ गइल रहनी, हमरा लगे इस्टेसन पर जातरी
लोग के आवाजाही देखे के फुरसत रहे। कुछ लोग \VUgras{असबाब-घर}
\textsuperscript{(25)} में आपन छोट सामान धरत भा उठावत रहे ; दोसर लोग
\VUgras{समय-सारनी} \textsuperscript{(26)} देखत रहे। आवे वाला जातरी लोग
\VUgras{निकास} \textsuperscript{(32)} की ओर जल्दी करत रहे, हाथ में
आपन \VUgras{सूटकेस} \textsuperscript{(24)} लेके। कुछ लोग \VUgras{असबाब बाँटे के
जगहा} \textsuperscript{(27)} पर अगोरत रहे आ आपन \VUgras{रसीद}
\textsuperscript{(29)} देके आपन संदूक, \VUgras{पेटी} \textsuperscript{(28)}
भा \VUgras{टोकरी} \textsuperscript{(30)} लेत रहे।

%%K t12-07-1 p
७. --- \VUgras{चुंगी के बाबू} \textsuperscript{(33)} लोग जतन से गठरी के
जाँचत रहे कि कुछ बतावे लायक बा कि ना।

%%K t12-08-1 p
८. --- कुछ जातरी \VUgras{जलपानघर} \textsuperscript{(34)} की ओर गइलें,
जहाँ एगो \VUgras{बैरा} \textsuperscript{(35)} ओह लोग के परोसे में लागल
रहे। ओह लोग के जल्दी करे के रहे, काहेकि \VUgras{गाड़ी}
\textsuperscript{(36)}, जवन तेज गाड़ी ह, इस्टेसन पर देर ले ना रुकी।

%%K t12-09-1 p
९. --- ओकरा बाद हम रवानगी के प्लेटफार्म पर गइनी आ सोझ जाए वाला
\VUgras{डिब्बा} \textsuperscript{(37)} खोजनी ताकि जातरा में डिब्बा बदले के
ना पड़े। हम एगो गलियारा वाला डिब्बा में चढ़नी जवना में पिए वाला लोग
खातिर एगो \VUgras{खाना} \textsuperscript{(38)} बा आ « अकेली मेहरारू »
खातिर एगो अलग खाना।

%%K t12-tit-2 sub
\VUpk{10.2pt}{40}{राति के रवानगी।}
\VUsaut{3.0mm}

%%K t12-10-1 p
१०. --- \VUgras{पायदान} \textsuperscript{(40)} पर चढ़ के, \VUgras{केवाड़}
\textsuperscript{(39)} बंद करके, \VUgras{परदा} \textsuperscript{(41)} लपेट
के आ आपन छोट असबाब \VUgras{जाल} \textsuperscript{(43)} पर धर के, हम
\VUgras{बेंच} \textsuperscript{(42)} पर बइठ गइनी। \VUgras{बत्ती वाला}
\textsuperscript{(44)} के बत्ती जरावत देख के हमरा लागल कि हमरा डिब्बा में
राति बितावे के बा, आ हम \VUgras{तकिया} \textsuperscript{(45)} किराया पर
देवे वाला बाबू के बोला के एगो तकिया माँगनी।

%%K t12-11-1 p
११. --- गाड़ी के कंडक्टर टिकट जाँचे आइलें। हम ओकरा से पूछनी कि हमनी के
अबहीं ले काहे नइखी चलत, काहेकि ठीक समय हो गइल बा। ऊ हमरा जवाब देलें कि
\VUgras{गार्ड} \textsuperscript{(46)} \VUgras{सामान के डिब्बा}
\textsuperscript{(47)} में साइकिल धरवावे में लागल बाड़ें। ऊपर से अबहीं
सब गोड़ गरम करे के डिब्बा \textsuperscript{(48)} बदलल नइखे।

%%K t12-12-1 p
१२. --- बाकिर हमनी के जल्दिए चले वाला रहनी जा, काहेकि \VUgras{आगवाला}
\textsuperscript{(50)} आपन \VUgras{टेंडर} \textsuperscript{(49)} पर से
कोयला उठा रहल रहे \VUgras{इंजन} \textsuperscript{(54)} के आग तेज करे
खातिर, आ \VUgras{ड्राइवर} \textsuperscript{(51)} खाली \VUgras{उप-इस्टेसन
मास्टर} \textsuperscript{(52)} के इसारा के अगोरत रहे, जे
\VUgras{प्लेटफार्म} \textsuperscript{(53)} पर रहलें।

%%K t12-13-1 p
१३. --- इंजन के \VUgras{बॉयलर} \textsuperscript{(56)} धीमे-धीमे गुर्रात
रहे ; \VUgras{चिमनी} \textsuperscript{(55)} \VUgras{भाप} \textsuperscript{(60)}
से मिलल धुँआ के धार उगलत रहे। एगो तेज \VUgras{सीटी} \textsuperscript{(59)}
गूँजल आ \VUgras{पिस्टन} \textsuperscript{(58)} चले लागल, ओह भारी मसीन के
\VUgras{पहिया} \textsuperscript{(57)} के ठेलत।

%%K t12-tit-3 sub
\VUsaut{3.0mm}
\VUpk{10.2pt}{40}{चल पड़लीं !}
\VUsaut{3.0mm}

%%K t12-14-1 p
१४. --- \VUgras{पटरी} \textsuperscript{(64)} पर दउड़त गाड़ी के रफ्तार तेज
हो गइल ; \VUgras{प्लेटफार्म} \textsuperscript{(65)} गायब हो गइल ; हमनी के
\VUgras{पाखाना} \textsuperscript{(66)} के पार कइनी जा, आ ओकरा बाद दोसरा
\VUgras{लाइन} \textsuperscript{(63)} पर खड़ा गाड़ी के भी : \VUgras{सुते के
डिब्बा} \textsuperscript{(67)}, \VUgras{खाए के डिब्बा}
\textsuperscript{(68)}, \VUgras{डाक के डिब्बा} \textsuperscript{(69)}।

%%K t12-noto-1 noto
\VUnotes{143.2mm}{%
\textit{(*) नोट।} --- \textit{समझ लीं कि जातरा के ई बयान लड़ाई से
पहिले के सरहद के हिसाब से बा।}%
}

%%K t12-15-1 p
१५. --- ओकरा बाद \VUgras{माल के इस्टेसन} \textsuperscript{(71)} हमनी के
आँख के सामने से गुजरल। छप्पर के नीचे बाबू लोग ऊ \VUgras{माल}
\textsuperscript{(72)} लादत रहे जवन धीमा गाड़ी से भेजल जाला।
\VUgras{मजूर के टोली} \textsuperscript{(76)} \VUgras{पानी के टंकी}
\textsuperscript{(73)} लगे \VUgras{जनावर के डिब्बा} \textsuperscript{(75)} आ
\VUgras{चपटा डिब्बा} \textsuperscript{(79)} के इधर-उधर करत रहे ताकि एगो
\VUgras{माल गाड़ी} \textsuperscript{(80)} बने।

%%K t12-16-1 p
१६. --- हमनी के आखिरी \VUgras{दुफाड़} \textsuperscript{(78)} पार कइनी जा,
जेकरा लगे काँटा वाला ठाढ़ रहलें ; हमनी के \VUgras{सिगनल}
\textsuperscript{(86)} के पार कइनी जा आ इस्टेसन दूर जाके गायब हो गइल।

%%K t12-17-1 p
१७. --- जल्दिए हमनी के ऊ \VUgras{लोहा के पुल} \textsuperscript{(74)} पार
कइनी जा जवन \VUgras{नदी} \textsuperscript{(81)} के ऊपर से जाला। ऊ पुल पार
कइला के बाद हमनी के नदी के साथे-साथे चलनी जा, जवना पर ताकतवर
\VUgras{खींचे वाला जहाज} \textsuperscript{(82)} भारी \VUgras{मालवाहक नाव}
\textsuperscript{(83)} खींचत रहे। हमनी के एगो \VUgras{जातरी जहाज}
\textsuperscript{(84)} भी देखनी जा, जब ऊ एगो \VUgras{जेटी}
\textsuperscript{(85)} लगे लागत रहे।

%%K t12-18-1 p
१८. --- कई गो इस्टेसन के तेजी से पार कइला के बाद हमनी के कुछ
\VUgras{सुरंग} \textsuperscript{(87)} पार कइनी जा आ आपन पाछे \VUgras{फाटक
के रखवार} के अनगिनत \VUgras{गुमटी} \textsuperscript{(88)} छोड़त गइनी जा।
गाड़ी के झोंका से झूलत हम गहिर नींद में सो गइनी।

%%K t12-tit-4 sub
\VUsaut{3.0mm}
\VUpk{10.2pt}{40}{दोयच-आवरिकूर इस्टेसन।}
\VUsaut{3.0mm}

%%K t12-19-1 p
१९. --- हम अचानक एगो बाबू के आवाज से जाग गइनी, जे चिचिआत रहलें :
« दोयच-आवरिकूर ! \textsuperscript{(*)}। सब लोग उतरीं ! » चुंगी के जाँच आ
सरहद के सब उबाऊ काम भइल। हमरा आपन जेब के घड़ी इस्टेसन के ओह
\VUgras{बड़का घड़ी} \textsuperscript{(89)} से मिलावे के पड़ल, जवन पचपन
मिनट आगे रहे ; नींद से मुसकिल से जागल हम \VUgras{बड़की सूई}
\textsuperscript{(90)} के \VUgras{छोटकी} \textsuperscript{(91)} से मुसकिल
से अलग कर पवनी ; भोर के कुहासा से \VUgras{डायल} \textsuperscript{(92)}
खुदे पूरा धुँधला रहे।

%%K t12-20-1 p
२०. --- जब चुंगी वाला लोग आपन जाँच करत रहे, हम जँभाई लेत ओह
\VUgras{इस्तिहार} \textsuperscript{(93)} के पढ़त रहनी जवन इस्टेसन के भीत
सजावत बा। समय काटे खातिर हम कुछ खा लेनी, जर्मन \VUgras{रेल के जाल}
\textsuperscript{(94)} के नकसा देखनी कि स्ट्रासबुर्ग अबहीं केतना दूर बा,
आ आपन डिब्बा में लवट गइनी, एह आस से ताकत पाके कि जल्दिए हम आपन जातरा
के मंजिल पर पहुँच जाइब।
